\section{Instance generation}
\label{app:instance}

This appendix formalises the accepted-instance family informally introduced in Section~\ref{sec:benchmark}. The full assembly procedure lives in \texttt{src/causal\_discovery/benchmark/instance.py}; the per-level configuration is fixed by the ladder in \texttt{run\_ladder.py}.

\subsection{DAG sampling}
\label{app:instance:dag}

Each instance starts from a random topological order $\sigma$ over $d$ node ids, drawn uniformly from $S_d$. The sampler enumerates the $\binom{d}{2}$ forward pairs $(i, j)$ with $\sigma^{-1}(i) < \sigma^{-1}(j)$ and selects exactly $k$ of them uniformly at random without replacement to form the edge set. The sampling is uniform on the set of DAGs with exactly $k$ forward edges under a uniformly random topological order; it is not Erd\H{o}s--R\'enyi.

The five in-paper levels fix $(d, k)$ as follows.

\begin{center}
\begin{tabular}{lcccccc}
\toprule
                        & L1 & L2 & L3 & L4  & L5  \\
\midrule
Nodes $d$              & 6  & 8  & 10 & 12  & 14  \\
Edges $k$              & 6  & 9  & 12 & 14  & 16  \\
Edge ratio $k/d$       & 1.00 & 1.13 & 1.20 & 1.17 & 1.14 \\
\bottomrule
\end{tabular}
\end{center}

\subsection{Linear--Gaussian SCM parameterisation}
\label{app:instance:scm}

For each edge $i \to j$ in the sampled DAG, the structural coefficient $\beta_{ji}$ is drawn uniformly from $[-2.0, -0.5] \cup [0.5, 2.0]$, with the sign chosen by an independent fair coin. Edges absent from the DAG receive coefficient zero. Each node receives an independent Gaussian noise term $\varepsilon_j \sim \mathcal{N}(0, \sigma^2_j)$ with intercept zero. For all in-paper levels, $\sigma^2_j = 1.0$ for every node (the $\sigma^2 = 0.5$ setting in the L0 spec is not reported).

The structural equations are
\begin{equation}
x_j = \sum_{i \in \mathrm{pa}(j)} \beta_{ji} \, x_i + \varepsilon_j,
\qquad j = 1, \ldots, d.
\end{equation}

\subsection{Rejection filters}
\label{app:instance:filters}

Three filters are applied in sequence. A candidate is rejected and resampled if any filter fires.

\textbf{Graph filter.} The CPDAG of the sampled DAG must contain at least two undirected edges, those undirected edges must not all share a single common node, and the DAG skeleton must be connected. This excludes instances whose observational equivalence class trivially identifies the DAG and instances whose skeleton splits into disjoint components.

\textbf{SCM filter.} The implied covariance matrix must be non-singular, and every \emph{relevant structural partial correlation} must satisfy $|\rho| \geq \varepsilon_{\mathrm{faith}}$ with $\varepsilon_{\mathrm{faith}} = 0.1$. The second condition is a finite-sample faithfulness margin: it removes parameter draws whose dependencies are too weak for any reasonable conditional-independence test to detect at the available sample sizes.

\textbf{Intervention-profile filter.} The minimum sufficient intervention set for the CPDAG (the lexicographically first set of single-node interventions that orient every undirected edge under Meek closure) must be non-empty. An empty intervention set means the observational CPDAG already pins down the full DAG, which collapses the active-discovery loop.

\subsection{Anonymisation and accepted-instance tuple}
\label{app:instance:anon}

A surviving candidate is anonymised by drawing a uniform random permutation $\pi \in S_d$ and relabelling the DAG, the CPDAG, the SCM, and the intervention set under $\pi$. The agent observes the public node ids only; the inverse $\pi^{-1}$ is retained by the evaluator and never exposed.

The accepted-instance tuple is
\begin{equation}
\mathcal{I} = \big(\, G_\pi,\; \mathcal{C}_\pi,\; \mathcal{M}_\pi,\; I^*_\pi,\; D_{\mathrm{obs}},\; B,\; \pi \,\big),
\end{equation}
where $G_\pi$ is the relabelled true DAG, $\mathcal{C}_\pi$ its CPDAG, $\mathcal{M}_\pi$ the relabelled linear--Gaussian SCM, $I^*_\pi$ the relabelled minimum sufficient intervention set, $D_{\mathrm{obs}} \in \mathbb{R}^{n_{\mathrm{obs}} \times d}$ a fresh observational sample drawn from $\mathcal{M}_\pi$, $B$ the per-episode intervention budget, and $\pi$ the anonymisation permutation. The runtime exposes $D_{\mathrm{obs}}$ and $B$ to the agent at episode start; everything else is hidden until scoring.

\subsection{Budget semantics}
\label{app:instance:budget}

The intervention budget is
\begin{equation}
B = |I^*_\pi| + s,
\end{equation}
where $s$ is the per-level slack. All in-paper levels use $s = 2$, so the agent always has at least two interventions in excess of the minimum sufficient set. Each intervention call to the runtime draws $n_{\mathrm{int}} = 25$ samples from $\mathcal{M}_\pi$ under $\mathrm{do}(V_j = v)$ and decrements the remaining budget by one.
